% Opcje klasy 'iithesis' opisane sa w komentarzach w pliku klasy. Za ich pomoca
% ustawia sie przede wszystkim jezyk i rodzaj (lic/inz/mgr) pracy, oraz czy na
% drugiej stronie pracy ma byc skladany wzor oswiadczenia o autorskim wykonaniu.
% \documentclass[declaration,shortabstract]{iithesis}
\documentclass[declaration,shortabstract,lic,english]{iithesis}

\usepackage{natbib}
\usepackage[utf8]{inputenc}

%%%%% DANE DO STRONY TYTUŁOWEJ
% Niezaleznie od jezyka pracy wybranego w opcjach klasy, tytul i streszczenie
% pracy nalezy podac zarowno w jezyku polskim, jak i angielskim.
% Pamietaj o madrym (zgodnym z logicznym rozbiorem zdania oraz estetyka) recznym
% zlamaniu wierszy w temacie pracy, zwlaszcza tego w jezyku pracy. Uzyj do tego
% polecenia \fmlinebreak.

\polishtitle    {Nienadzorowane uczenie przez dopasowywanie\fmlinebreak rozkładów prawdopodobieństw}
\englishtitle   {Unsupervised learning using output distribution matching}
\polishabstract {W nieniejszej pracy sprawdzamy hipotezę, czy można nauczyć sieć neuronową rozpoznawać tekst bez korzystania z podpisanych próbek tekstu, a jedynie z rokładów prawdopodobieństw wystąpienia pewnych sekwencji liter (n-gramów) w danych uczących. Implementujemy rozwiązanie zaproponowane przez  Yu Liu, Jianshu Chen oraz Li  Deng  w pracy  “Unsupervised  Sequence Classification using Sequential Output Statistics” oraz testujemy je na zbiorze danych MNIST, który jest zbiorem danych o dużo mniejszej regularności niż zbiór danych używany przez w cytowanej pracy. Wkładamy również wysiłek w próbę zrozumienia zasad działania algorytmu przez wiele różnych przypadków testowych (innych od przedstawionych w cytowanej pracy) oraz przez użycie bardziej skomplikowanego modelu sieci neuronowej.
}

\englishabstract{This thesis validates a hypothesis whether it is possible to teach a neural network optical character recognition without labeled data, instead using sequential output statistics (in the form of n-grams) only. We implement solution proposed by  Yu Liu, Jianshu Chen and Li Deng in their paper “Unsupervised Sequence Classification using Sequential Output Statistics” and we test it on on MNIST data set, which is much less regular data set than the one used by Liu et al. [2017]. We also put effort into understanding algorithm's effects through multiple different test cases (different than those in cited paper) and we try to use more complex neural net model.}

\author         {Grzegorz Ciesielski}
% w przypadku kilku promotorow, lub koniecznosci podania ich afiliacji, linie
% w ponizszym poleceniu mozna zlamac poleceniem \fmlinebreak
\advisor        {dr Jan Chorowski}
\date          {}                     % Data zlozenia pracy

\transcriptnum {290956}
\advisorgen    {dr. Jana Chorowskiego}
%%%%%

%%%%% WLASNE DODATKOWE PAKIETY
%
%\usepackage{graphicx,listings,amsmath,amssymb,amsthm,amsfonts,tikz}
%
\usepackage[]{algorithm2e}
\usepackage{graphicx,amsmath,subcaption,multirow}

%%%%% WŁASNE DEFINICJE I POLECENIA

\newcommand{\calD}{\mathcal{D}}
\makeatletter
\newcommand{\pushright}[1]{\ifmeasuring@#1\else\omit\hfill$\displaystyle#1$\fi\ignorespaces}
\newcommand{\pushleft}[1]{\ifmeasuring@#1\else\omit$\displaystyle#1$\hfill\fi\ignorespaces}
\makeatother

\begin{document}

%%%%% POCZĄTEK ZASADNICZEGO TEKSTU PRACY

\chapter{Introduction}

Most machine learning models have two main requirements: a well designed model and (usually a lot of) data. Preparing the model usually does not cause any problems, as there are a lot of them available online and their performance allows them to learn on semi-advanced, affordable machines. But to do so, we also need the second requirement: data. Especially in optical character recognition (OCR) it seems that we cannot build any decent model without some labeled samples, and gathering such data requires a lot of human effort. Yet there is some property of written text which we could exploit -- it is very regular and, from a probabilistic standpoint, quite sparse because it is difficult to randomly sample a sequence of letters which will create a correct English word (or sequence of words making a reasonable sentence). If we had known that, for example, recognised text is written in Spanish, we could take Spanish dictionary (which is easily available) and penalize our model for predicting words that do not appear in the dictionary, and award it for the opposite. Or if we had a lot of text in Spanish, we could also compute its n-gram statistics and use them to supervise model. This approach was proposed by Yu Liu, Jianshu Chen and Li Deng in paper ``Unsupervised Sequence Classification using Sequential Output Statistics'' \citep{liu2017unsupervised}. In this thesis we will study this possibility of training neural classifier to recognize images of written text without any labeled data, but only using sequential statistics regarding output. We will implement and test solution proposed by~\citet{liu2017unsupervised}, except that we will use smaller but more diverse data set and more complex models.

\section{Problem formulation}

We consider the problem of learning an optical character recognizer using data consisting of unlabeled sequences of data and the desired output statistics. In language processing obtaining such data set is considerably simpler that labeled data, as it is enough to get images of written text in specific language and for output statistics use easily obtainable n-grams of this language. Specifically: our classifier predicts sequences \( (y_1, \dots y_n) \) (labels) from input sequences \( (x_1, \dots x_n) \). The  algorithm we are going to implement will have access to:
\begin{enumerate}
    \item a data set \( \mathcal{D} = \lbrace (x_1^k, \dots x_n^k) : k = 1, \dots M \rbrace \) of sequences,
    \item n-gram probabilities of such data - which are denoted as: 
\[ p(i_1, \dots i_n) = p(y_1^k = i_1, \dots y_n^k = i_n) \]
where \( i_1, \dots i_n \in \lbrace 0, 1, \dots C \rbrace \) and \( y_j^k \) is the label of letter \( x_j^k \).
\end{enumerate}

In the OCR task, $x_j^k$ corresponds to images of letters, $y_j^k$ to actual labels, and $p(i_1, \dots i_n)$ to the probability of the sequence $i_1, \dots i_n$ occurring in valid English sentence.

Similarly to \citet{liu2017unsupervised} we search for a sequence classifier in the form of $p_\theta(y_t^k | x_t^k)$ which computes posterior probability based only on input sample $x_t^k$. What will be different than \citet{liu2017unsupervised} is that they used linear classifiers, whereas we use neural network with 2 hidden layers (\mbox{LeNet-300-100} architecture \citep{lecun1998gradient}) and a data set containing more diversified samples.

\chapter{Problem solution}


Based on given data, what first comes to mind is to create a model that minimizes negative cross entropy between prior distribution and expected n-gram frequencies:
\begin{equation} \label{cross-entropy}
    \mathcal{J}(\theta) := - \sum_{i_1, \dots i_n}p(i_1,\dots i_n) \ln \overline{p}_\theta(i_1, \dots i_n)
\end{equation}

\citet{liu2017unsupervised} have shown that the expected n-gram frequency $\overline{p}_\theta(i_1, \dots i_n)$ can be described in terms of $p_\theta(y_t^k = i_k | x_t^k)$ as:

\begin{equation} \label{expected-probability}
    \overline{p}_\theta(i_1, \dots i_n) = \frac{1}{M}\sum_{k=1}^{M}\prod_{j=1}^{n}p_\theta(y_j^k = i_j | x_j^k)
\end{equation}

Now it is possible to evaluate loss function $\mathcal{J}(\theta)$ (Equation \ref{cross-entropy}), but we encounter significant problem during evaluation of its derivative - the sample average is inside logarithm. This makes it incompatible with stochastic gradient descent. Furthermore \citet{liu2017unsupervised} noticed that $\mathcal{J}(\theta)$ is highly non-convex and there are high barriers between local optima and global minimum. Proposed solution uses Legendre transformation \citep{nielsen2010legendre, boyd2004convex} which, for convex function $f(u)$ is defined as:
\begin{equation} \label{legendre-transformation}
    f^*(v) := \sup_{u}(v^Tu - f(u))
\end{equation}
Taking $f(u)=-\ln u$ gives us $f^*(v) = -1 - \ln(-v)$. It holds that $(f^*)^* = f$ \citep{boyd2004convex}, hence:
\begin{equation} \label{logarithm-transform}
    - \ln u = f(u) = \sup_v(uv - f^*(v)) = \max_v(uv + 1 + \ln(-v))
\end{equation}

We can now rewrite (\ref{cross-entropy}) by substituting logarithm into equation (\ref{logarithm-transform}), which transforms our minimization problem into min-max problem:

\begin{equation} \label{min-max-definition}
    \min_\theta \max_{\lbrace v_{i_1, \dots , i_n} < 0 \rbrace} \left\lbrace 
    \mathcal{L}(\theta, V) := \frac{1}{M} \sum_{k=1}^M L_k(\theta, V) + \sum_{i_1, \dots, i_n} p(i_1, \dots, i_n) \ln(-v_{i_1, \dots, i_n})
    \right\rbrace 
\end{equation}

where $V := {v_{i_1, \dots i_n}}$ are dual variables and $L_k(\theta, V)$ is the loss function of $k$-th sequence:
\begin{equation}
    L_k(\theta, V) := \sum_{i_1,\dots,i_n} p(i_1,\dots,i_n) \cdot v_{i_1,\dots,i_n} \cdot \prod_{j=1}^{n} p_\theta(y^k_j = i_j | x^k_j)
\end{equation}

Now we have no summation under logarithm. Moreover, the new dual loss function $L(\theta, V)$ was shown by \citet{liu2017unsupervised} to be smoother and more suitable for stochastic gradient descent algorithm, but we have to adapt SGD to working with two loss functions: primal $\mathcal{L}_V^{primal}(\theta) := \mathcal{L}(\theta, V)$ and dual $\mathcal{L}_\theta^{dual}(V) := \mathcal{L}(\theta, V)$ -- the first being maximized, the second minimized. \citet{liu2017unsupervised} proposed an algorithm described below:

\begin{algorithm} 
 \KwData{$\calD = \lbrace (x_1^k, \dots x_n^k) : k = 1, \dots M \rbrace$ and $p(i_1,\dots i_n)$}
 Initialize $\theta$ and $V$ \\
 \While{not converged}{
 $\mathcal{B} \leftarrow \lbrace x^{k_i}_1, \dots x^{k_i}_n \rbrace _{i=1}^B$ ($B$ sequences
 sampled randomly from data set $\calD$) \\
 Compute average of gradients from minibatch:
 $$\Delta\theta = \frac{1}{B} \sum_{i=1}^{B} \frac{\partial L_{k_i}}{\partial\theta}, \hspace{0.1cm}
 \Delta V = \frac{1}{B} \sum_{i=1}^{B} \frac{\partial L_{k_i}}{\partial V} + \frac{\partial}{\partial V} \sum_{i_1,\dots i_n} p(i_1, \dots i_n) \ln(-v_{i_1, \dots i_n})
 $$ \\
 $\theta \leftarrow \theta - \mu_\theta \Delta \theta$ \\
 $V \leftarrow V + \mu_V\Delta V$
 }
 \caption{Stochastic Primal-Dual Gradient Descent}
\label{spdg}
\end{algorithm}


\chapter{Tests}

\section{Experiment details}

In the tests we implemented the Stochastic Primal-Dual Gradient Descent (SPDG) algorithm (\ref{spdg}) in Python 3.7 with PyTorch 1.1. Our model (primal variable) is a neural net with 2 hidden layers, \mbox{LeNet-300-100} architecture \citep{lecun1998gradient}, which is different that architecture used in \citet{liu2017unsupervised}, because they used only linear classifiers and mentioned that algorithm will not be suitable for more complex architectures. We worked with MNIST data set \citep{mnist} - grayscale $28 \times 28$ images of digits, transformed into $n \times 28 \times 28$ sequences (Fig. \ref{fig:mnist_sequences_example}) and gathered into minibatches containing 128 sequences each (during tests on Brown Corpus in Section~\ref{section:real_life} we increased this size to 1024). We used ADAM optimizer for both primal and dual variables, with learning rates $\mu_\theta=10^{-6}$ and $\mu_V = 10^{-4}$ respectively. We initialized dual variables with uniformly distributed values from interval $(-1, 0)$.

Due to algorithms unpredictability and heavy dependence on initialization of variables (mainly holding for dual variables), we decided to show only single runs (not average of some trials) on diagrams, as they were very difficult to average - but we will note whether other runs were similar or not. All experiments and diagrams can be viewed at our GitHub repository. We will reference them as ``GH a.b.c'', where ``a.b.c'' is the name of the folder containing the test.

\begin{figure}[htb]
    \centering
    \includegraphics{doc/images/img_1537094603.png}
    \includegraphics{doc/images/img_8385974286.png}
    \includegraphics{doc/images/img_9227163033.png}
    \caption{Example sequences of MNIST digits.}
    \label{fig:mnist_sequences_example}
\end{figure}

\section{Unigrams}

\begin{figure}[htb]
    \def\svgwidth{\columnwidth}
    \input{images/t112_test_error.pdf_tex}
    \caption{Test error on unigrams (GH t1.1.2a). In this test case, data set contained only zeros and ones, with proportions 9:1 noted in n-gram. This is only one run of algorithm, but other runs yielded very similar results (but difficult to average).}
    \label{fig:t112_error}
\end{figure}


\citet{liu2017unsupervised} mentioned that SPDG algorithm would not work properly if applied to 1-grams (which would be equivalent to deciphering Caesar's Cipher). We would like to put some effort into understanding why and how it does not work, especially because when applied it produces very peculiar results, as shown in Figure \ref{fig:t112_error}.

We can note from the diagram that the model does not seem to learn at all from the beginning, it stays right about the same for 500 epochs, and then it `crashes', causing errors to change drastically during one epoch every ~20 epochs or so. What happens? Let us take a closer look at our loss function (simplified for our case):

\newcommand{\sumxinB}{\sum_{x \in \mathcal{B}}}
\newcommand{\bpmean}{p_\theta^\mathcal{B}}
\begin{align} 
\mathcal{L}^\mathcal{B}(\theta, V) &= \frac{1}{B}\left(\sumxinB p(0)\cdot v_0\cdot p_{\theta}(y=0 | x) + \sumxinB p(1)\cdot v_1\cdot p_{\theta}(y=1 | x)\right)\\ 
& \pushright{+p(0)\cdot\ln(-v_0) + p(1)\cdot\ln(-v_1)=} \\
&= p(0) \cdot\left(\bpmean(0)\cdot v_0+\ln(-v_0)\right) + p(1)\cdot\left(\bpmean(1) \cdot v_1+\ln(-v_1) \right)
\end{align}

Here (and below) $\bpmean(i) = \frac{1}{B}\sumxinB p_\theta(y=i | x)$ is the mean posterior probability of $i$.

In the algorithm we used primal learning rate of $\mu_\theta=10^{-6}$, while dual learning rate was 100 times higher at $\mu_V = 10^{-4}$ - it was for dual variables to adjust to the posterior probability distribution. Thus at the beginning we can assume that $\bpmean$ is roughly constant through all minibatches. Hence our algorithm starts to maximize $\mathcal{L}(\theta, V)$ with respect to dual variables $V$ (where every $v_i \in V$ is from interval $(-\infty, 0)$). We can note that:

\begin{align}
    \lim_{v_i\rightarrow 0} \mathcal{L}(\theta, V) &= \lim_{v_i \rightarrow 0} (p_\theta(i) \cdot v_i + \ln (- v_i)) = -\infty \\
    \lim_{v_i\rightarrow-\infty} \mathcal{L}(\theta, V) &= \lim_{v_i \rightarrow -\infty} (p_\theta(i) \cdot v_i + \ln (- v_i)) = -\infty
\end{align}

Hence its maximum is where the gradient nullifies. Differentiating over $v_i$:

\newcommand{\sumxinD}{\sum_{x \in \mathcal{D}}}
\begin{align}
    \frac{\partial\mathcal{L}}{\partial v_i}(\theta, V) = \frac{\partial}{\partial v_i} (p(i) \cdot (p_\theta(i) \cdot v_i + \ln(-v_i))) = 
    p(i) \cdot (p_\theta(i) + \frac{1}{v_i})
\end{align}

\begin{figure}[tb]
    \def\svgwidth{\columnwidth}
    \input{images/t112_dual.pdf_tex}
    \caption{Dual variables in the test case with unigrams, where data set consisted of only zeroes and ones (GH t1.1.2a, the same as in Figure \ref{fig:t112_error}).}
    \label{fig:t112_dual}
\end{figure}


It is equal zero at $(v_0, v_1) = (-p_\theta(0)^{-1}, -p_\theta(1)^{-1})$. As we see in the Figure \ref{fig:t112_dual}, optimal $v_0$ is reached after circa 50 epochs. After that we noticed that model stopped learning zeroes completely, it was only fortifying his current knowledge -- adjusting primal variables to output higher values approaching infinity and causing numerical problems with softmax and gradient evaluation, which resulted in this strange behaviour of large error jumps when variables overflow.

\paragraph{Conclusions.} In the case of unigrams SPDG algorithm will not achieve decent accuracy. Its behaviour can be described as follows:
\begin{enumerate}
    \item Sample primal and dual variables, producing some posterior (random) distribution;
    \item Adjust dual variables to match this distribution;
    \item Increase confidence to infinity.
\end{enumerate}

The numerical overflow could easily be mitigated by putting some constraints on model's parameters or simple detection whether model produces high enough values right before applying softmax, but this would not improve accuracy.

\section{3-grams}

For most of our tests we decided to use trigrams, which is due to the fact that summation in equation \ref{min-max-definition} goes over every n-gram entry. Thus because the n-gram size increases exponentially with $n$, we are motivated to use smaller $n$ to decrease the complexity of computations. Hence because bigrams were too simple for most of our tests, trigrams became the perfect choice.

MNIST database consists of 60000 training images (ca. 6000 for each digit) and 10000 testing images \citep{mnist}. That gives us $2.16\cdot10^{11}$ different possible assignments of images to just one 3-digit sequence, linearly increasing with n-gram size (number of entries) and exponentially with $n$. Obviously, we used only a subset of those (sampled randomly), but this creates an issue. If, for example, we've created a data set of 50000 3-digit sequences using (which is already 2.5 as much data as MNIST data set) by random choice with replacement of digits, every digit would occur roughly 15000 times (for some of them this number might be significantly smaller). Given that we have 6000 images of this number and we sample 15000 of them with replacement, we will use only 5500 of them (as on average 500 samples will not be chosen). If there would be a digit occurring 10000 times only, then it would use only $\sim81\%$ of all its images available.

\subsection{Simple 3-gram}

During this test case (GH t3.1.1, Fig. \ref{fig:t311_errors}) training data set $\mathcal{D}$ contained 40000 3-digit sequences, from which 36000 were $(0, 1, 2)$, and 4000 were $(1, 2, 3)$ giving 9:1 ratio noted in n-gram: $p(0,1,2)=0.9$, $p(1, 2, 3)=0.1$ and $p(\overline{i})=0$ otherwise.

\begin{figure}[htb]
\begin{subfigure}[b]{.49\textwidth}
    \def\svgwidth{\textwidth}
    \input{images/t311_test_error.pdf_tex}
    \caption{Error rate on test set.}
    \label{fig:t311_test_error}
\end{subfigure}
\begin{subfigure}[b]{.49\textwidth}
    \def\svgwidth{\textwidth}
    \input{images/t311_train_error.pdf_tex}
    \caption{Error rate on data set (input sequences).}
    \label{fig:t311_data_error}
\end{subfigure}
\caption{Error rate on test and data set for test case (GH t3.1.1) with digits 0-3 and 3-gram containing 2 nonzero entries: $p(0,1,2)=0.9$ and $p(1, 2, 3)=0.1$.}
\label{fig:t311_errors}
\end{figure}

As we can see in Figure \ref{fig:t311_test_error}, test set errors for 0, 1 and 2 reach decent errors almost immediately, while the number 3 requires about $100$ epochs to achieve test error below $10\%$. To answer the question why this happens, let's differentiate loss function by $\theta$:

\newcommand{\sumxsinB}{\sum_{(x1, x2, x3)\in \mathcal{B}}}
\begin{align}
    \frac{\partial\mathcal{L}}{\partial \theta}(\theta, V) &= \frac{\partial}{\partial \theta} \left(p(0, 1, 2) \cdot v_{(0, 1, 2)} \cdot \bpmean(0, 1, 2) + p(1, 2, 3)\cdot v_{(1, 2, 3)} \cdot \bpmean(1, 2, 3)\right) \\
    &= \frac{9}{10} \cdot v_{(0, 1, 2)} \cdot \frac{\partial \bpmean}{\partial \theta}(0, 1, 2) + \frac{1}{10}\cdot v_{(1, 2, 3)} \cdot \frac{\partial \bpmean}{\partial \theta}(1, 2, 3)
\end{align}

Because $v_i < 0$, minimizing loss means maximizing posterior probabilities $\bpmean(0,1,2)$ and $\bpmean(1,2,3)$ with respect to $\theta$ with weights $\frac{9}{10} \cdot v_{(0, 1, 2)}$ and $\frac{1}{10}\cdot v_{(1, 2, 3)}$ respectively. Given that $\frac{9}{10} \cdot v_{(0, 1, 2)} > \frac{1}{10}\cdot v_{(1, 2, 3)}$ (which happens $94\%$\footnote{Let's say that $X, Y \sim \mathcal{U}(0,1)$. Then $P\left(\frac{9}{10}X > \frac{1}{10}Y\right) = P\left(9X > Y\right) = \frac{8}{9}P\left(9X > Y | X > \frac{1}{9}\right) + \frac{1}{9}P\left(9X > Y | X \le \frac{1}{9}\right) = \frac{8}{9} + \frac{1}{9}\cdot\frac{1}{2} = 0,9444\dots$}
of the time if we sample dual variables from uniform distribution from interval $(-1, 0)$), our model will aim towards predicting sequence $(0,1,2)$ every time. We can see that in Figure \ref{fig:t311_data_error}, at first error rate on data set for $3$ is equal $100\%$, which confirms our assumptions. 

\begin{figure}[htb]
\begin{subfigure}[b]{.49\textwidth}
    \def\svgwidth{\textwidth}
    \input{images/t311_test_error.pdf_tex}
    \caption{Error rate on test set (GH t3.1.1).}
    \label{fig:t311_test_error_comp}
\end{subfigure}
\begin{subfigure}[b]{.49\textwidth}
    \def\svgwidth{\textwidth}
    \input{images/t322_test_error.pdf_tex}
    \caption{Error rate on test set (GH t3.2.2).}
    \label{fig:t322_test_error_comp}
\end{subfigure}
\caption{Comparison between test errors for two test cases (GH t3.1.1 and GH t3.2.2) with only two nonzero n-gram entries $(0,1,2)$ and $(1,2,3)$. The main difference between them was that in former $p(0,1,2) = 0.9$ and $p(1,2,3) = 0.1$, while in latter $p(0,1,2) = 0.6$ and $p(1,2,3) = 0.4$}
\label{fig:t311_t322_test_error_comp}
\end{figure}

What would happen, if our model would start learning less frequent sequence from the beginning? We have conducted a similar tests but with $p(0,1,2)=60\%$ and $p(1,2,3)=40\%$. Here probability of swapping distributions at the beginning was $33\%$\footnote{As previously, $P(\frac{3}{5}X>\frac{2}{5}Y) = P(X>\frac{2}{3}) + P(\frac{3}{2}X>Y|X\le\frac{2}{3})=\frac{1}{3} + \frac{1}{2}\cdot\frac{2}{3} = \frac{2}{3}$.} When model learns correctly, it learns quicker and better than its counterpart with higher ratio - e.g. in Figure \ref{fig:t311_t322_test_error_comp}.

\begin{figure}[htb]
\begin{subfigure}[b]{.49\textwidth}
    \def\svgwidth{\textwidth}
    \input{images/t342_train_error.pdf_tex}
    \caption{Error rate on train set.}
    \label{fig:t342_data_error}
\end{subfigure}
\begin{subfigure}[b]{.49\textwidth}
    \def\svgwidth{\textwidth}
    \input{images/t342_test_error.pdf_tex}
    \caption{Error rate on test set.}
    \label{fig:t342_test_error}
\end{subfigure}
\begin{subfigure}[b]{.49\textwidth}
    \def\svgwidth{\textwidth}
    \input{images/t342_predictions_train_012.pdf_tex}
    \caption{How is every (0,1,2)-sequence classified by model. Output sequences shown only if they occurred a total
    of at least $12000$ times across all epochs.}
    \label{fig:t342_predictions_012}
\end{subfigure}
\begin{subfigure}[b]{.49\textwidth}
    \def\svgwidth{\textwidth}
    \input{images/t342_predictions_train_123.pdf_tex}
    \caption{How is every (1,2,3)-sequence classified by model. Output sequences shown only if they occurred a total
    of at least $10000$ times across all epochs.}
    \label{fig:t342_predictions_123}
\end{subfigure}
\caption{Results from test case (GH t3.2.4) with digits 0-3, 3-gram containing 2 nonzero entries: $p(0,1,2)=0.6$ and $p(1, 2, 3)=0.4$. Here model does not learn correctly - mismatches sequences to distributions.}
\label{fig:t342_results}
\end{figure}

But when the model starts by learning the less frequent sequence first, it will most likely try to predict that sequence $(0, 1, 2)$ (which is 1.5 times more frequent than $(1, 2, 3)$) is actually $(1, 2, 3)$. We experienced such result during one test case which results are shown in Figure \ref{fig:t342_results} (GH t3.4.2). Let us note from diagram \ref{fig:t342_predictions_012} that indeed from the very beginning almost every sequence $(0,1,2)$ was classified as $(1, 2, 3)$, resulting in model classifying $(1, 2, 3)$ as $(2, 3, x)$, where $x \in \lbrace0, 1, 2, 3\rbrace$ - which figure \ref{fig:t342_predictions_123} confirms. 

\paragraph{Overfitting} In figures \ref{fig:t311_test_error_comp}, \ref{fig:t322_test_error_comp} and \ref{fig:t311_test_error} we can notice that after roughly 700 epochs test errors for 1 and 2 significantly increase. As error on train set is constantly very low, such results are caused by model overfitting to train data set. Should it concern us? Actually, not as much as in other machine learning problems. For us here, test error is only a determinant of regularization of our model. Very valuable is our ability to label the train set almost perfectly, which we could have then used for supervised learning.

\paragraph{Interesting fact.} In figure \ref{fig:t342_predictions_123} we can see that no $(1, 2, 3)$ was classified as $(2, 3, 0)$ almost ever (it was predicted a total of 113 times, only during first 10 epochs and never again - whereas $(2, 3, 1)$ was predicted a total of a little over 700.000 times). In fact, model quite quickly completely `forgot' to predict zeroes, as it was unable to classify anything as $(0, 1, 2)$ due to the fact that every sequence seen by model was either $(1, 2, 3)$ or started with $(2, 3, \dots)$, giving no mapping to $(0, 1, 2)$ possible - which resulted in zeroes being redundant.


\subsection{Cyclic 3-grams}

\begin{table}[htbp]
\centering
\begin{tabular}{c|c}
    Sequence & Frequency \\
    \hline
    (0, 1, 2) & 20\% \\
    (1, 2, 3) & 20\% \\
    (2, 3, 4) & 20\% \\
    (3, 4, 0) & 20\% \\
    (4, 0, 1) & 20\% \\
\end{tabular}
\caption{Sequence distribution during test cases GH t$3.5.x$, for $x=1\dots$.}
\label{tab:cyclic_test}
\end{table}

For some test cases (GH t$3.5.x$, for $x=1\dots$) we generated data set $\mathcal{D}$ containing 40000 3-digit sequences with digits 0-4. Data set was generated with sequence frequencies described in Table \ref{tab:cyclic_test}.

In 4 out of 5 of runs, errors approached 100\% quickly and maintained such for duration of training. When evaluated, model has shown that for each digit it predicts mainly one, different digit (both on data and test). The assignment of actual digit to predicted by model was, in every case, some cyclic permutation of the sequence (0, 1, 2, 3, 4). If we counted posterior frequencies of sequences, it would exactly match input distribution (Table \ref{tab:cyclic_test}), and regarding that our model does not know that it `swapped' some digits it will never recover from this situation. From its perspective it matched distributions perfectly, as it reached global optimum of loss function, whereas from ours sight of view the results were completely incorrect.

In one run model successfully learnt to recognize letters, with minimum mean error $2.09\%$ and minimum greatest error $2.77\%$.

\subsection{Cyclic 3-gram with broken symmetry of distributions}

Similarly to the previous test, we tried to train model with cyclic 3-gram, but with one difference in sequence frequencies: one entry was 1.5 times more likely to occur than any other (tests GH t$3.6.x$, Table \ref{tab:cyclic_test_no_symmetry}). Results of this run are shown in Figure \ref{fig:t362_results}. It learnt to correctly predict zeroes, ones and twos (which were part of the most frequent sequence), and than based on this knowledge matched distributions of remaining sequences and learnt to predict correctly other digits. Minimum mean and minimum max error were higher than in previous test (at $2,72\%$ and $4,17\%$ respectively), but that was due to more imbalanced distribution of letters -- instead of every letter occurring roughly same amount, 3s and 4s occured half as much as in previous tests, hence resulting in higher errors on 3s and 4s. Errors on remaining digits were lower, but not as much to compensate.

\begin{table}[htbp]
\centering
\begin{tabular}{c|c}
    Sequence & Frequency \\
    \hline
    (0, 1, 2) & 60\% \\
    (1, 2, 3) & 10\% \\
    (2, 3, 4) & 10\% \\
    (3, 4, 0) & 10\% \\
    (4, 0, 1) & 10\% \\
\end{tabular}
\caption{Sequence distribution during test cases GH t$3.6.x$, for $x=1\dots$.}
\label{tab:cyclic_test_no_symmetry}
\end{table}

\begin{figure}[ht]
\begin{subfigure}[b]{.49\textwidth}
    \def\svgwidth{\textwidth}
    \input{images/t362_train_error.pdf_tex}
    \caption{Error rate on train set.}
    \label{fig:t362_train_error}
\end{subfigure}
\begin{subfigure}[b]{.49\textwidth}
    \def\svgwidth{\textwidth}
    \input{images/t362_test_error.pdf_tex}
    \caption{Error rate on test set.}
    \label{fig:t362_test_error}
\end{subfigure}
\caption{Results of test case with cyclic 3-gram entries, but broken symmetry in distributions (Tab. \ref{tab:cyclic_test_no_symmetry}, GH t3.6.2).}
\label{fig:t362_results}
\end{figure}


\section{Real life examples}
\label{section:real_life}

We wanted to test our model also on more lifelike examples, specifically English language. We decided to use Brown Corpus \citep{francis1964brown} from Natural Language Toolkit for Python \citep{loper2002nltk}, which is a data set of English text consisting of 1.100.000 words. From those we produced a sequence of 4.750.000 characters. We filtered $C$ most commonly occurring characters and mapped them to digits $0, 1, \dots C-1$, giving us sequence of digits whose n-gram distribution simulates distribution of (some) letters in English language. As sequences produced consisted of at least 1.5 million up to 3.5 million characters (depending on the number of filtered characters $C=3,4\dots10$), we decided to use fraction of it as train set (in tests it was the first 100000 characters) and some other fraction as test set (characters 100000--110000 in our tests). 

\paragraph{N-gram choice.} Due to selecting only a fraction of Brown Corpus, we were to decide how to generate n-gram probabilities. We could have used the whole Brown Corpus and evaluate its n-gram, or generate n-gram probability distribution only on chosen subset of data. First approach would be more accurate with real life example, as when we learn sequence classifier on unlabeled data (some English text) we would not have access to this data's n-gram probabilities - we would have had only English n-gram available. Second approach is also valuable - it shows maximum abilities of SPDG algorithm (when it has access to perfect distribution).

\subsection{N-gram generated from the entire data set}

In the following tests (GH: tbs3.3.1, tbs3.4.1 and tbs3.5.1) we fed model with 3-grams evaluated on whole Brown Corpus generated sequence of digits, while during training we mapped only the first 100000 characters to MNIST digits. Each test case was working with different number of digits: 3, 4 and 5 respectively ($C=3, 4, 5$). We did not manage to teach neural net with higher dimension of output ($C > 5$).

\begin{table}[h]
\centering
\begin{tabular}{c|c|c|c|c|c|c|}
    \multicolumn{2}{c}{} & \multicolumn{5}{|c|}{Distribution (\%)} \\
    \cline{3-7}
    Test case & \% of available sequence & 0 & 1 & 2 & 3 & 4 \\
    \hline
    BS3.3.1 & 7.07 & 41.4 & 30.3 & 28.2 & - & - \\
    BS3.4.1 & 5.63 & 33.1 & 24.2 & 22.6 & 20.1 & - \\
    BS3.5.1 & 4.71 & 27.8 & 20.4 & 19.0 & 16.9 & 15.9 \\
\end{tabular}
\caption{Distribution of characters (digits) in tests on Brown Corpus.}
\label{tab:tbs3(345)1_stats}
\end{table}

Table \ref{tab:tbs3(345)1_stats} below shows more accurate statistics about train set for each test case. During training model has had access to whole MNIST train set of every digit, as there are $\sim6000$ images of every, while the least we used were fives in test case BS3.5.1, which was 15.9\% of 100000-digit sequence - that gives 15900 images used, still more than 2.5 times as much as available images (hence some were repeated).


\begin{figure}[ht]
\begin{subfigure}[b]{.49\textwidth}
    \def\svgwidth{\textwidth}
    \input{images/tbs331_test_error.pdf_tex}
    \caption{Error rate on test set during test case with number of digits $C=3$ (GH tbs3.3.1.)}
    \label{fig:tbs331_test_error}
\end{subfigure}
\begin{subfigure}[b]{.49\textwidth}
    \def\svgwidth{\textwidth}
    \input{images/tbs341_test_error.pdf_tex}
    \caption{Error rate on test set during test case with number of digits $C=4$ (GH tbs3.4.1.)}
    \label{fig:tbs341_test_error}
\end{subfigure}
\begin{subfigure}[b]{.49\textwidth}
    \def\svgwidth{\textwidth}
    \input{images/tbs351_test_error.pdf_tex}
    \caption{Error rate on test set during test case with number of digits $C=5$ (GH tbs3.5.1.)}
    \label{fig:tbs351_test_error}
\end{subfigure}
\begin{subfigure}[b]{.49\textwidth}
    \def\svgwidth{\textwidth}
    \input{images/tbs3(345)1_mean_test_error.pdf_tex}
    \caption{Comparison of mean error rates of previous tests.}
    \label{fig:tbs3(345)1_mean_test_error}
\end{subfigure}
\caption{Results of using SPDG algorithm to predict sequences of MNIST images mapped on Brown Corpus. N-grams used here were generated using the whole Corpus, but data used in training was only a subsequence of the whole Corpus.}
\label{fig:tbs1_results}
\end{figure}

In Figure \ref{fig:tbs1_results} we see results of 3 tests on Brown Corpus based MNIST sequences:
\begin{enumerate}
\item $C=3$ case (GH tbs3.3.1, Figure \ref{fig:tbs331_test_error}): results were replicable and net usually achieved minimum mean test error on level $\sim6\%$ after 40-80 epochs and then balanced around $6-8\%$. Error was diversified - for zeroes and ones it was below $5\%$, while twos never achieved error rate lower than $10\%$ in any run.
\item $C=4$ case (GH tbs3.4.1, Figure \ref{fig:tbs341_test_error}): here model was more dependant on initialization. During testing it learnt to below $20\%$ average error roughly $50\%$ of the time (2 out of 4 trials), while other $50\%$ it stayed far above $50\%$ and only learnt to predict one letter. Here minimum mean error was at higher level of $8\%$ after 60-100 epochs, and then maintained below $12\%$.
\item $C=5$ case (GH tbs3.5.1, Figure \ref{fig:tbs351_test_error}): during testing we managed to get a working example only once. Results of this run are shown on diagrams: lowest mean error was $\sim 10\%$ after 80 epochs and then stayed below $15\%$ for entire training.
\end{enumerate}

Also worth noting are mean train errors for every test case: $3.6\%, 5.25\%$ and $6.2\%$ respectively, achieved after 50-100 epochs and remaining roughly constant during entire training. And why is that train error is interesting for us? This is because we used unlabeled data during training, hence our ability to label it with over $90\%$ accuracy might be very profitable for us. We could have, for example, created program that without any previous knowledge (except for maybe n-gram distributions of English language) would be able to recognize images of text, given that we get enough images to match n-gram distribution in some extent. Of course, we omit obvious problem of segmentation of data, but such algorithms already exist (e.g. Tesseract from \citet{kay2007tesseract}).

\subsection{N-gram generated from train set}

\begin{figure}[ht]
\begin{subfigure}[b]{.49\textwidth}
    \def\svgwidth{\textwidth}
    \input{images/tbs362_test_error.pdf_tex}
    \caption{Error rate on test set during test case BS3.6.2. (number of digits $C=6$)}
    \label{fig:tbs362_test_error}
\end{subfigure}
\begin{subfigure}[b]{.49\textwidth}
    \def\svgwidth{\textwidth}
    \input{images/tbs382_test_error.pdf_tex}
    \caption{Error rate on test set during test case BS3.8.2. (number of digits $C=8$)}
    \label{fig:tbs382_test_error}
\end{subfigure}
\caption{Results of SPDG algorithm on MNIST images mapped on Brown Corpus, but with access to exact train data n-gram distribution.}
\label{fig:tbs2_results}
\end{figure}

To test limits of SPDG algorithm on Brown Corpus, we decided to give model access to the most accurate n-gram statistics possible -- the one generated directly from train set. It increased model's consistency significantly. Tests with $C=3,4,5$ were not as often crashing (GH tbs3.3.2, tbs3.4.2, tbs3.5.2), and a bit more accurate (usually 1-2 percentage points on mean error). We even managed to train sequence classifier with $C=6$ and $C=8$ up to reasonable accuracy - model with $C=6$ (GH tbs3.6.2, Fig. \ref{fig:tbs362_test_error}) achieved minimum mean error at $8.25\%$ (almost 2 percentage points better than test case with $C=5$, but with less accurate 3-gram), while model with $C=8$ learnt to predict 7 out of 8 letters with errors below $30\%$ (GH BS3.8.2, Fig. \ref{fig:tbs382_test_error}).

\chapter{Conclusions}

In our thesis we tested the \citet{liu2017unsupervised} method on nonlinear classifier classifier (here neural net with 2 hidden layers) with unlabeled data with successful results. Even though we did not fully simulate classification of sequences of English text (we only managed to classify 5-7 letters correctly), we have shown other properties of algorithm with its ability to train complex classifiers. Although MNIST might seem as simple data set, it is actually quite difficult to reduce test errors below $7\%$ (linear classifiers did not -- \cite{mnist}). We achieved low error rates even though we complexified prediction problem through removing access to labels, but having only (expected) output sequence probability distributions.

Furthermore, most of our models have shown that test and train error rates were very similar. This proves that our model generalized quite well and did not overfit too much (especially during real life training tests). Furthermore, some MNIST images are extremely similar yet with different labels (even for humans). This gives us hope that when applied to images of more consistently written text, SPDG algorithm will achieve even better results.

\bibliographystyle{plainnat}
\bibliography{lic}

\end{document}
